\section{19. The donor stack as a transformation family}

V3's composition support is a \texttt{5 x 5} loop over donor families in which, in the
committed checker, neither donor appears on the right-hand side of anything and
both legs are the same loop-invariant expression. Instrumented, that loop
evaluates \texttt{compose} at \textbf{2 of its 8 possible argument triples}. Under the
interpretation each family becomes a transformation with its own hand-off
contracts, and the third argument is computed from the pair rather than typed.

Three frame conditions are stated as axioms rather than folded into an encoding:

\begin{enumerate}
  \item \textbf{Hand-offs are never contract identities} --- no family's target contract is
    any family's source contract, its own included.
  \item \textbf{Distinct families have distinct endpoints.}
  \item \textbf{The stack is inhabited.}
\end{enumerate}

\subsection{Theorem V4.4 --- the composition rows, generalised}

A bridged hand-off between two carrying donors composes; an unbridged one
refuses; a donor composed with itself still needs a bridge registered from its
own target contract back to its own source contract; and a chain of three donors
carries exactly when all three legs carry and both interior hand-offs are
bridged. For a stack of any size.

\subsection{Theorem V4.5 --- the rows the enumeration does not contain}

If either leg fails to carry, the composite does not carry, whatever the other
leg is and whatever the registry says. These are the six argument triples the
published loop never evaluates.

\subsection{Theorem V4.6 --- the donor axis indexes something}

Distinct ordered donor pairs present distinct hand-offs, and two pairs whose
hand-offs the registry decides differently have composites that carry
differently. On an inhabited carrying stack, both the successes and the refusals
have witnesses rather than holding vacuously.

Two conditions that looked like frame conditions --- that no identity is a donor,
and that a composite's hand-off is separated too --- were written as axioms first,
reported inert, and are now discharged as theorems from the remaining three.
Their inertness is re-measured rather than remembered: both are added back to the
axiom set on every run and the theorems that becomes newly provable are reported,
and the answer has to stay none.
